% Three doors into one column of data, and a guard standing in front of
% exactly one of them.
%
% The Speakers subgraph exposes three fields on Query, all three advertised
% by introspection, and all three reach the same rows of the Speaker table:
%   speakers(first:)                the root field
%   _entities(representations:)     the entity route the router uses
%   node(id:)                       the node field
% A guard placed on the root field only stops speakers(first:) and does
% nothing to the other two: _entities and node still return Speaker.email in
% full. That is the whole picture: one path blocked, two paths open, one
% column of data behind all three.
%
% Layout is boxes and links, the idiom figures/tikz/ch11-position-is-the-join.tex
% uses, not the axis-and-ticks timeline of ch01-one-field.tex and
% ch06-two-routes.tex: the subject here is three parallel routes converging
% on one target, not a sequence unfolding over one request, so there is no
% single shared axis for all three to sit on.
%
% Three field boxes on the left, one tall box on the right standing for the
% Speaker rows and their email column (drawn tall on purpose, so a single
% box can receive an arrow at more than one height and still read as one
% target). The entities and node links run straight through from field box
% to data box, arrowhead at the data box: the request is answered. The
% speakers link runs only as far as a small guard box and stops there: it
% never reaches the data box at all, which is the point a reader should
% get in one glance, not from a label.
%
% A short dashed mark crosses the entities and node links at the same
% x-coordinate as the guard box, purely as an annotation (the links
% themselves are not broken there). It marks where a check would also have
% to sit if it were checking the data rather than one field, and the note
% beside it says so. This is the one place this figure borrows the "guard
% belongs around the data, not on one path" idea the brief allows; it stays
% inside one panel rather than adding a second, because the same guard box
% and the same gate x-coordinate already carry the comparison.
%
% No fills, square corners, one stroke weight throughout (0.4pt), arrowhead
% only where a request actually lands (the data box, or the guard box for
% the one path that gets no further) -- the same "arrowhead marks the one
% directional claim" idiom ch11 uses. Dashed is a short, local mark here,
% not a link of its own, and does not carry the same meaning dashed carries
% in ch11; the note beside it is what supplies the meaning, same as ch11's
% own legend does for its two link styles.
%
% No versions, timings or counts anywhere, per SPEC decision 19.
%
% No quotation marks anywhere in this file (strict Ascii mode; a literal "
% reads as a quote character to the prose gate, and \" is a diaeresis
% accent that fails to compile inside a node). \_entities is written with
% the escape TikZ needs for an underscore in text mode, the same way
% ch06-two-routes.tex escapes \_\_typename.
%
% Widths: field boxes are 38mm text plus 2mm inner sep, centered at x=22,
% so they span x=1 to x=43, starting near the left margin the way every
% other figure in this book does. The guard box is centered at x=84 and
% spans x=78 to x=90. The data box is 26mm text plus 2mm inner sep,
% centered at x=135, spanning x=120 to x=150 -- the same right-hand edge
% ch11's note uses, and inside the 155mm text measure (A4, 30mm inner,
% 25mm outer). The note beside the dashed marks runs from x=90 to x=136,
% clear of that edge.
%
% Bare tikzpicture (house style): the figure environment, caption and
% label live at the call site in the section file.
\begin{tikzpicture}[
  x=1mm, y=1mm,
  every node/.append style={font=\sffamily\scriptsize},
  lane/.style={anchor=west, font=\sffamily\scriptsize\itshape},
  fieldbox/.style={draw, line width=0.4pt, align=center, text width=38mm,
    inner sep=2mm, minimum height=14mm},
  gatebox/.style={draw, line width=0.4pt, align=center, text width=10mm,
    inner sep=1mm, minimum height=10mm},
  databox/.style={draw, line width=0.4pt, align=center, text width=26mm,
    inner sep=2mm, minimum height=100mm},
  reachlink/.style={draw, line width=0.4pt,
    -{Straight Barb[length=1.4mm, width=1.6mm]}},
  blocklink/.style={draw, line width=0.4pt,
    -{Straight Barb[length=1.4mm, width=1.6mm]}},
  gapmark/.style={draw, line width=0.4pt, dashed},
  note/.style={align=left, inner sep=2pt},
]

% ---------------------------------------------------------------- heading
\node[lane] at (0,58) {Speakers subgraph, three fields of Query};

% ------------------------------------------------------------ field boxes
\node[fieldbox] (fspk)  at (22,40)  {speakers(first:)\\the root field};
\node[fieldbox] (fent)  at (22,0)
  {\_entities(representations:)\\the entity route};
\node[fieldbox] (fnode) at (22,-40) {node(id:)\\the node field};

% -------------------------------------------------------------- the guard
\node[gatebox] (guard) at (84,40) {guard};
\node[note, anchor=north] at (84,32) {refused here};

% ------------------------------------------------------------- data column
\node[databox] (data) at (135,0) {Speaker rows\\email column};

% ------------------------------------------------------------------ links
% The two open routes: unbroken, arrowhead at the data box.
\draw[reachlink] (fent.east)  -- (120,0);
\draw[reachlink] (fnode.east) -- (120,-40);

% The guarded route: stops at the guard box and goes no further.
\draw[blocklink] (fspk.east) -- (guard.west);

% ------------------------------------------------------- where a check on
% ------------------------------------------------- the data would also sit
\draw[gapmark] (84,-5)  -- (84,5);
\draw[gapmark] (84,-45) -- (84,-35);
% text width 26mm, not 46: the note is anchored at the gate x and the data
% box starts at x=120, so anything wider than 30mm runs into the box.
\node[note, anchor=west, text width=26mm] at (90,-20)
  {the same crossing point on both open routes, and nothing stands there};

\end{tikzpicture}
